\subsection{ROS 2 Ring Buffer}
\label{sec:ros-ring-buffer}
\kedu{TODO this section is not final. This section lists the functions we need to verify for the concurrent ring buffer, if we decide to choose it as a case study. Probably start with \lstinline|enqueue| and \lstinline|dequeue|, followed by \lstinline|ctor| and others if time permitting. }

We use \kedu{$Paper Name$} to prove functional correctness of in-production C++ code: the Robot Operating System 2 (ROS 2)'s concurrent ring
buffer.
\footnote{\url{https://github.com/ros2/rclcpp/blob/a525a2439d0fd91041f3a125d549e6abf8a0a46f/rclcpp/include/rclcpp/experimental/buffers/ring_buffer_implementation.hpp\#L38-L238}} 
RAII-style mutexes protects concurrent accesses to the buffer's internal state, permitting multiple threads to perform \lstinline|enqueue| and \lstinline|dequeue| on a shared ring buffer. \kedu{Maybe also mention non-concurrent features like overflow?}

\begin{lstlisting}[
  language=C++,
  basicstyle=\scriptsize\ttfamily,
  aboveskip=.5\baselineskip,
  belowskip=.5\baselineskip,
  caption={Thread-safe operations of the ROS 2 ring buffer.},
  label={lst:ros-ring-buffer}
]
template<typename BufferT>
class RingBufferImplementation :
  public BufferImplementationBase<BufferT> {
public:
  explicit RingBufferImplementation(size_t capacity)
  : capacity_(capacity),
    ring_buffer_(capacity),
    write_index_(capacity_ - 1),
    read_index_(0), size_(0) {
    if (capacity == 0)
      throw std::invalid_argument(
        "capacity must be a positive, non-zero value");
  }
  virtual ~RingBufferImplementation() {}

  void enqueue(BufferT request) override {
    std::lock_guard<std::mutex> lock(mutex_);
    write_index_ = next_(write_index_);
    ring_buffer_[write_index_] = std::move(request);
    if (is_full_())
      read_index_ = next_(read_index_);
    else
      size_++;
  }

  BufferT dequeue() override {
    std::lock_guard<std::mutex> lock(mutex_);
    if (!has_data_())
      return BufferT();
    auto request = std::move(ring_buffer_[read_index_]);
    read_index_ = next_(read_index_);
    size_--;
    return request;
  }

  std::vector<BufferT> get_all_data() override {
    // return copy of data[read_index_, read_index_+size_]
  }

  inline size_t next(size_t val) {
    std::lock_guard<std::mutex> lock(mutex_);
    return next_(val);
  }
  inline bool has_data() const override {
    std::lock_guard<std::mutex> lock(mutex_);
    return has_data_();
  }
  inline bool is_full() const {
    std::lock_guard<std::mutex> lock(mutex_);
    return is_full_();
  }
  size_t available_capacity() const override {
    std::lock_guard<std::mutex> lock(mutex_);
    return available_capacity_();
  }
  void clear() override {
    std::lock_guard<std::mutex> lock(mutex_);
    clear_();
  }

private:
  inline size_t next_(size_t val) {
    return (val + 1) % capacity_;
  }
  inline bool has_data_() const {
    return size_ != 0;
  }
  inline bool is_full_() const {
    return size_ == capacity_;
  }
  inline size_t available_capacity_() const {
    return capacity_ - size_;
  }
  inline void clear_() {
    size_ = 0;
    read_index_ = 0;
    write_index_ = capacity_ - 1;
  }
  // Representation fields are omitted.
};
\end{lstlisting}
